\section{Baseline Difficulty and Candidate Advantage}

The near-zero full-roster effect of the comoving ego predictor
combines different effects on easy and hard rows. In the same
OpenDD development folds, we define easy and hard using CV errors
and cuts estimated on the other sites. They are outcome strata,
not inference features. Contributions are normalized by the
whole site's observed CV error mass, so they add rather than
reporting percentages on different denominators.

\begin{table}[h]
\centering
\caption{Additive lower gain contributions of the frozen neural
forecast versus CV. Three seeds are averaged before the
seven-site bootstrap.}
\small
\begin{tabular}{lrl}
\toprule
Stratum & Rows & Contribution [95\% interval] \\
\midrule
Easy & \DifficultyEasyRows{} & \DifficultyEasyEffect{} \\
Hard & \DifficultyHardRows{} & \DifficultyHardEffect{} \\
Other observed & \DifficultyOtherRows{} & \DifficultyOtherEffect{} \\
Wholly unobserved & \DifficultyUnobservedRows{} & \DifficultyUnobservedEffect{} \\
\bottomrule
\end{tabular}
\end{table}

A sign guard splits the residual along and perpendicular to
the last observed motion and keeps each component if it agrees
with the corresponding OLS4-minus-CV displacement. Its direct
lower effect versus the original neural direction at exactly
the same rowwise peaks is \DifficultyDirection{}.
At the \DifficultyEdits{} fixed slots it reaches
\DifficultyGuard{} joint passes versus \DifficultyPrefix{} for
the donor-prefix rule, with direct effect \DifficultyGuardMatched{}.
The difference in counts does not establish improved utility.

We separately train \DifficultyFits{} past-only heads to predict CV
difficulty, not neural gain. The head obtains hard AUROC
\DifficultyAUROC{} versus \DifficultyMotionAUC{} for normalized
CV/OLS4 disagreement. Using the neural forecast when predicted hard
probability is at least 0.50 and easy probability below 0.25 yields
\DifficultyRawEdits{} switches. Against the same trajectory bank
and switch count per query ranked by disagreement, its direct
gain is \DifficultyGatedMotion{}; versus CV it is
\DifficultyGatedCV{}. At the original radii and slots it changes
\DifficultySwaps{} shapes, obtaining \DifficultyGate{} versus
\DifficultyPrefix{} joint passes in \DifficultyViews{} views
and a direct effect of \DifficultyGatedFixed{}.
These controls distinguish predicting baseline error from
predicting a specific replacement's advantage. Neither rule
advances to external calibration.
